% chapters/ch13.tex —— 第 13 章
% 体例：具体失败 → 设计问题 → 机制解释 → 对照实验 → 条件变化 → 选择原则
% 数字身份纪律：外部数字一律 \cited{}；本书实验台的实测用 \measured{}
% （来源 experiments/results/triggers-E01-20260923T152230Z-90acda.json，
%   原始记录 experiments/runs/E01-20260923T152230Z-90acda/）。
% E11 尚未运行，本章不给出任何关于 P5 的判决。
\chapter{何时进化：触发条件与资源预算}
\label{ch13}

\chaptergoal{每次遇到新经验，都值得更新系统吗？}
\chapterprops{P5关注不同任务流中怎样选择更新时机。本章给出每任务、周期与事件触发三种时机的成本形状和可计算信号；E11尚无真实模型比较。}
\chapterfig{图~\ref{fig:ch13-trigger}：三种更新时机在时间轴上的落点，以及触发判决要回答的两个问题}
\chapterpitfall{八次写入里没有一次被新颖性门拦下，于是“每条都该写”看起来是对的}

% ---------- 1. 具体失败 ----------
\begin{failurebox}
每条新记忆都立即触发旧记忆改写，可能使写入成本随记忆存储增长。D-MEM 用 $O(N^2)$ 描述同步持续演化的规模瓶颈，同时把一次深演化写为 $O(N)$；本书将二者分别放在累计增长与单次操作的层次理解，不能合成“每次写入都需要平方时间”的结论。其分层路由方案在指定带噪配置下报告约 \cited{80\%}的接口 token 降幅\cite{dmem-2603}。实际增长率仍取决于每次链接、检索和改写的范围；首次写入也有构建与调用费用。

还需要核对筛选规则实际区分了什么。
本书对一次真实运行的训练轨迹作了事后回算：
E01 运行 \texttt{E01-20260923T152230Z-90acda} 的格 4 在实际执行的八道训练题上保存了 \measured{8} 条经验记录，
按字符二元组重叠算，这 \measured{8} 条的新颖性最小 \measured{0.8209}、均值 \measured{0.9546}。
用这个信号做 Gate（低于 0.5 就不写），\textbf{一条都不会被拦下}。

这不证明这些经验在语义上都新颖，也不证明该信号在任何阈值下均无用；它只说明这批文本在所选阈值下全部被放行。\textbf{原运行没有部署这个门，回算不能称为实际节省了多少调用或改善了多少任务。}
\end{failurebox}

\designquestion{更新时机需要由可执行规则定义，规则中可以包含周期、阈值等超参数。它同时决定成本账的形状与收益账的形状，
因此需要明确可执行的调度规则；采用事件触发时，尤其要说明\textbf{运行期能得到哪些信号}以及它们何时可用。}

\section{三种时机：每任务、周期与事件触发}
\label{sec:ch13-timing}

\begin{table}[htbp]
\centering
\caption{三种更新时机：成本长什么样、什么时候失效}
\label{tab:ch13-timing}
\small\setlength{\tabcolsep}{3pt}
\begin{tabular}{>{\raggedright\arraybackslash}p{1.7cm}>{\raggedright\arraybackslash}p{3.0cm}>{\raggedright\arraybackslash}p{3.1cm}>{\raggedright\arraybackslash}p{3.1cm}>{\raggedright\arraybackslash}p{3.2cm}}
\toprule
时机 & 本章操作化 & 成本形状 & 失效条件 & 需要的账 \\
\midrule
每任务 & E11 每个任务后请求构建 & 请求数随任务数增长；实际费用依准入与额度 & 若无去重，重复任务可能产生重复条目 & 写入次数 / 去重后条数 \\
周期 & E11 每隔固定题数请求构建 & 请求数随周期数增长；每次成本可变 & 周期过长时可能延迟响应 & 周期内累计失败数 \\
事件触发 & E11 当前信号达到或超过阈值才请求 & 深更新费随触发次数增长，监测费仍随任务数增长 & 信号无分辨率时退化成每任务 & 触发率 / 漏触发数 \\
\bottomrule
\end{tabular}
\end{table}

三种时机不是三档优劣，而是三种把成本挂在哪个量上的选择。
本章每任务与周期分别按任务序号和固定题数提出请求；周期也可以按墙钟定义，但那是另一种调度。事件策略按信号越界提出请求。请求频率不等于实际发布频率或总费用。

相关论文提供了不同的时间边界，不能直接装进这三行。Reflexion在任务试次间产生反思，并不等于每个跨任务episode后都发布新经验\cite{reflexion-2303}。AgeMem在第一、二阶段之间重置上下文、保留长时记忆，阶段内允许工具操作；这种训练轨迹结构不能单独证明固定周期更新\cite{agemem-2601}。A-MAC研究候选准入，D-MEM研究分层路由，均需与本章完整调度策略区分。
触发次数减少只是可能节省深更新费用；若信号计算、失败重试或后续推理更贵，总账仍可能增加。触发率必须进账，并与实际资源和任务质量一起阅读。

事件触发有一种需要单独诊断的失效方式：\textbf{信号没有响应本应处理的变化}。每任务调度在无预算与执行故障时会提出更新请求，但请求仍可能被拒绝、失败或产生无用候选；周期调度也可能在下一周期到来前错过短暂机会。因此不能把其他策略视为绝不会遗漏。阈值过高的事件触发器则可能整段不提出请求，外部日志只看到“没有更新”，需要进一步区分稳定与失明。

\begin{sloppypar}
漏触发还有一个对称的兄弟：\textbf{误触发}——信号越过了阈值，但并没有新东西值得写。
它比漏触发更难发现，因为它的账看起来是正常的：
一次写入发生了、验证也跑了、持久物也更新了，
如果没有显式的重复与有效性检查，“这次写入是否带来新信息”就可能被漏记；文字去重也只能检查表面重复。
因此事件触发这一栏要同时报两个率：触发率与\emph{触发后的去重率}。
只报第一个时，“学会了挑时机”与“整体变懒了”仍然无法区分，
但“挑对了时机”与“触发了但没新东西”同样无法区分。
\end{sloppypar}

\subsection{E11 协议：这三种时机要怎么比}
\label{sec:ch13-e11}

\begin{sloppypar}
\emph{信号能算}不等于\emph{三种时机比过}。
本书第 \ref{ch04} 章之前的触发信号分析（\path{experiments/lab/trigger_signals.py}）
是在\textbf{已跑完的轨迹}上算出来的描述统计量；
E11 的旧协议（\path{experiments/protocols/E11.md}，2026-09-28 写成独立协议）
要求的四件事在此前都不具备：
\end{sloppypar}

\begin{itemize}
  \item \textbf{配对任务流}：比较单元是一条固定顺序的\emph{任务序列}，不是单个 episode；
        重采样单元必须是序列，下钻到 episode 会把相关样本当成独立样本。
  \item \textbf{资源约束}：旧协议规定共同写入次数上限 $W$，但相同写入上限不等于相同总调用、token 或时间预算。修订版须将做题、经验构建和监测的资源范围分开列明，并对所选择的总资源量施加可执行上限。
  \item \textbf{监测开销入账}：已有字符运算不调用模型，但仍有本地耗时；若计算新颖性以前必须先生成候选经验，该生成调用也须入账；
        若改用模型或向量检索算新颖性，必须按实际调用数入账。
  \item \textbf{漂移前后交互}：序列中点之后整体换族，判决输入是两段\emph{差值之差}，
        不是“漂移后谁赢”。
\end{itemize}

\begin{sloppypar}
协议把 P5 拆成两个可分别被推翻的子判决：
指定范围内的单位资源收益差，以及该差值在漂移前后是否扩大；完整排序还需要事件与周期、周期与每任务的相邻比较。
\textbf{E11 本身尚未运行}，因此本章不给出任何一种时机的实测优劣；
第 \ref{ch08} 章已报告 E05 v2 工程试跑，但仍没有足够证据给出技能固化的实测回本点——
协议就位不等于结果就位。
\end{sloppypar}


\subsection{让信号对应当前时点：一次运行器审计}

E11 的旧实现提供了一个具体的工程反例。运行器在每题结束后，把截至当前的判定列表与已经写入的经验列表交给同一个信号函数。前者一题一项，后者只有写入时才增长。函数逐项配对取最大值时，输出长度跟着较短的经验列表结束。如果已经完成多道题却只存过一条经验，返回的末项仍是最早那一项，而不是当前题的信号。程序没有报错，数组也不为空，事件臂却在读取错误时点。

这不是模型性能问题，也不能靠扩大样本弥补。一次不含模型调用的构造检查即可复现：完成\calcnum{4}题、只存\calcnum{1}条经验时，旧函数只返回\calcnum{1}个信号。第一条经验的新颖性按初始化规则为\calcnum{1.0}，于是取末项的调用者会反复读到这个旧值。这样的“事件触发”可能很快耗尽写入额度，却不能解释为系统发现了连续的新事件。

修复时首先要定义每个信号的观察对象。连续失败以已经结束的任务为索引；漂移窗口以判定时间序列为索引；经验新颖性以一个已经存在的候选经验为输入。若候选尚未生成，就不存在可免费读取的候选新颖性。可以先生成候选再决定发布，但必须给所有实际生成记费用；也可以用当前任务描述与历史任务描述计算输入新颖性，但它检验的是另一种信号，需要改名、预定阈值并记录设计变化。用已存经验的末项填补尚未生成的候选，会混淆时间与对象。

第二个错误发生在统计阶段。设每次运行前半段两臂都通过，后半段基线全失败而事件臂全通过。按本章失败率差的定义，每次运行的优势扩大为\calcnum{1}。若把这段序列重复两次后直接串接，再从整条列表中点切开，“前半段”变成完整第一次运行，“后半段”变成完整第二次运行，两边表现相同，结果被算成\calcnum{0}。正确方法是在每次运行内部识别切换点，先算前后差，再在预定配对单位内汇总重复运行。增加重复次数不能改变漂移点的含义。

第三个边界是更新时间对后续任务的作用。任务结束后形成的经验，最早影响下一道任务；序列最后一题之后的更新不可能改善已经完成的成绩。如果把末尾写入计作“有用更新”，或者用整段结束时才满足的漂移窗口解释前面的收益，就颠倒了时间顺序。末尾是否仍为未来部署保留更新，应由成本范围决定：只评价这段序列时可预先关闭；若要保存，则照常计费但不能给当前分数加收益。

离线审计还发现重复运行使用了相同的轨迹输出目录。原始轨迹必须包含序列、臂与重复编号，否则后一次可能覆盖前一次，汇总分数即使保留，也失去逐次核查入口。这些问题及输入示例保存在\nolinkurl{calculations/e11-readiness-audit-2026-09-30.json}。旧实现和协议保留为历史版本；修订控制器、分析与运行器必须完整验证并冻结之后，才可开始新的真实试跑。本节没有产生 P5 的能力证据。

\section{三个可计算的触发信号}
\label{sec:ch13-signals}

本章把三个触发信号落到脚本
\path{experiments/lab/trigger_signals.py} 上，跑在已有运行的真实轨迹里。
按任务顺序重建这些描述量时，不需要未来题目的判定；但候选新颖性要等经验已经生成才能计算，判定也必须已对相应角色可见。它们不是在任意时刻都可免费获取：

\begin{enumerate}
  \item \textbf{错误集中}：写入时往前数，末尾连续失败的次数。它回答“这次失败是不是孤立的”。
  \item \textbf{分布漂移}：最近 4 个 episode 的失败率与前 4 个的差的绝对值。它回答“失败率是否发生变化”；绝对值同时响应改善和恶化，不能单凭它判定失败变多。
  \item \textbf{经验新颖性}：1 减去本次经验与此前已有经验的最大字符二元组相似度。它回答文本与既有文本有多相似，不直接判断语义重复或有用性。
\end{enumerate}

\begin{table}[htbp]
\centering
\caption{三个信号在 E01 一次真实训练上的取值（格 4，\measured{8} 条已存记录）}
\label{tab:ch13-signals}
\begin{tabular}{>{\raggedright\arraybackslash}p{3.6cm}cccc}
\toprule
信号 & 最小 & 均值 & 最大 & 回算未放行比例 \\
\midrule
错误集中（连续失败数） & \measured{0} & \measured{0.3750} & \measured{2} & \multirow{3}{*}{\measured{0.0000}} \\
漂移（含暖机零值） & \measured{0.0000} & \measured{0.0625} & \measured{0.5000} & \\
经验新颖性 & \measured{0.8209} & \measured{0.9546} & \measured{1.0000} & \\
\bottomrule
\end{tabular}

\begin{minipage}{0.92\linewidth}
\vspace{2mm}
{\footnotesize 表注：原始记录 \path{experiments/runs/E01-20260923T152230Z-90acda/}，
渲染结果 \path{experiments/results/triggers-E01-20260923T152230Z-90acda.md}。
Gate 规则为事后示例（错误集中 $\ge 2$ 或新颖性 $\ge 0.5$ 才放行），不是实际部署协议；最后一列为对已存文本的回算。漂移行包含七个暖机零占位；只有第八题形成有效窗口比较，其绝对差为0.5，方向为失败减少。
样本只有 \measured{8} 次写入，只够说明信号的量级，不够支持任何判决。}
\end{minipage}
\end{table}

这张表最后一列表示：按该阈值回算，\calcnum{8}条记录全部放行，未放行比例为\calcnum{0}。数值集中在较高区间，因而\calcnum{0.5}阈值没有筛选作用；这同时涉及信号分布与阈值，不能说与阈值无关。用词改变可以降低表面重叠，但没有逐条语义标签，不能仅据这张表认定八条都是同义经验或信号完全失效。

由此得到一条可以写进协议的要求：
\textbf{任何声称用“新颖性”做触发的系统，必须同时报出新颖性的分布}。
只报“我们用了新颖性门”等于什么都没说——
本书这次只是固定历史输入上的回算，未评价Gate改变未来记忆后会产生的任务轨迹。

\subsection{暖机零值不是七次稳定观察}

旧脚本需要最近四题与此前四题两个完整窗口。当前记录只有\measured{8}题，因此前\calcnum{7}个位置没有满足比较条件；程序返回零作为暖机占位，只有最后一个位置产生实际窗口差。将八个数组元素相加再除以八得到\measured{0.0625}，这个数忠实复现了旧程序，却不能解释成八次可用漂移测量的均值。

原始判定依次为“通过、失败、失败、通过、通过、通过、通过、通过”。前窗失败率为\calcnum{2/4}，后窗为\calcnum{0/4}。有效绝对差是\calcnum{0.5}，带方向的后窗减前窗为\calcnum{$-0.5$}：这一次变化对应失败减少，不能称为失败突然增多。窗口不足、观测到零变化和反馈缺失应分别标记；新版报告可以同时保留历史算法读数与按有效窗口计算的量，不能覆盖旧输出冒充原来的分析。

八条经验的来源任务、序号与八条训练轨迹一一对应，没有触发旧脚本的任务号补位路径。若别的记录缺少经验文本，该脚本会用任务号补齐；这种输出不应继续命名为经验新颖性。公开预览与完整档案都收录了这次 E01 运行：由\path{analysis/inputs.json}找到对应的\path{summary.json}，其所在目录下的\path{traces/grid4/train/}保留八条训练轨迹。从仓库根目录运行以下命令，脚本会对齐经验与轨迹，并将复算结果写入\path{analysis/timing-audit-reproduction.json}。
\begin{verbatim}
python calculations/reproduce_timing_audit.py
\end{verbatim}
完整档案还保留早期复核报告\path{calculations/ch13-trigger-replay-audit-2026-09-30-1815.json}。

\subsection{信号会越过阈值，不等于环境真的变化}
\label{sec:ch13-signal-noise}

第\ref{sec:ch13-signals}节给出可计算的量，本节讨论这些量的解释边界。这里的概率均来自本书明确设定的独立伯努利教学模型，计算文件为\nolinkurl{calculations/ch13-calculations-2026-09-30.json}，不是 E11 的实测错误率。

\subsubsection{短窗口为什么容易产生大幅变化}

考虑一个完全稳定的任务流，每道题独立失败的概率都为 $p$。用最近两题失败比例减去此前两题失败比例，再取绝对值。每个窗口的比例只能是 $0$、$1/2$、$1$。本节规定差值 $D\ge1/2$ 时触发，等号也算达到阈值。只要两个窗口失败题数不同，就会触发；“比例相差二分之一”在这里仅意味着某个窗口多错了一题，而不是积累了很多证据。

当 $p=1/2$ 时，每个窗口的失败题数 $K$ 服从两次伯努利之和，$K=0,1,2$ 的概率分别是 $1/4,1/2,1/4$。两个互不重叠窗口相同的概率为
\[
 \Pr(K_1=K_2)=(1/4)^2+(1/2)^2+(1/4)^2=3/8.
\]
因此稳定环境中一次窗口比较达到阈值的概率为 $1-3/8=\calcnum{5/8}$，即\calcnum{62.5\%}。这是教学设定下的精确概率，说明大幅的短窗比例变化本身未必稀有；它不是对 E11 后端误报率的估计，也不是要求所有实际系统放弃这个信号。若将比较符改成严格的 $D>1/2$，只有一窗全失败、另一窗全通过才会触发，概率变为\calcnum{1/8}。所以窗口长度、阈值数值与是否包含等号，必须一起写进规则。

把稳定失败率改成 $p=1/4$，按同样的含等号规则有限枚举得到触发概率\calcnum{69/128}，约\calcnum{53.91\%}。基准失败率变低并不会在这个阈值和窗口长度下自动消除频繁触发。实际任务可能有相关性，滑动窗口也彼此重叠，这些条件会进一步改变连续报警的结构；不能把每次比较当作独立试验，用简单次数相乘推算整段误报。

“连续两次失败”也需要区分单个位置与整个运行期。在同一独立教学模型、$p=1/4$ 下，指定相邻两题都失败的概率是 $1/16$。但在\calcnum{8}题里至少出现一对相邻失败的概率，经枚举所有判定序列为\calcnum{20905/65536}，约\calcnum{31.90\%}。两个量的事件范围不同，把前者当作整段报警概率，会低估多次观察带来的触发机会。

这些计算不要求上线时照搬一个统计检验。它们要求阈值设计回答一个朴素问题：在我愿意称作“没有变化”的情形里，这条规则本来会响多少次？若答案不清楚，就无法区分触发器敏感与触发器过于频繁。可以在不用于最终评价的稳定任务序列上先画分布，也可以用明确的生成模型做压力测试；两者都要标明来源，不能把构造世界的概率当部署保证。

\subsubsection{输入变化、错误变化与经验变化}

本书旧轨迹上的“经验新颖性”比较候选经验文本，新版 E11 的“输入新颖性”比较当前公开任务描述。名字相近不意味着证据等价。任务描述换了表述，输入新颖性可能很高，而正确解法完全相同；两个任务描述几乎一样，边界条件却可能改变，需要不同做法。候选经验的文字差异也不能证明语义差异，更不能证明有用性。

把输入新颖性放在构建之前的优点，是不必为了决定是否构建而先支付构建费；代价是它远离最终要保留的经验内容。把候选新颖性放在构建之后，则可以比较实际候选，但节省的主要是后续验证、发布或存储费用，已经发生的生成费不会消失。实验若更换这两种做法，需要修改成本边界与信号定义，而不能仍称“同一个新颖性门”。

三个信号采用逻辑或时，还会出现主导信号。如果输入新颖性几乎每题都超过阈值，连续错误与窗口变化即使记录得很准确，也很少改变最终决策。此时发现“所有信号都启用了”不足以证明三者都有效。应逐题保留触发理由，计算仅由某信号触发的次数、多个信号同时成立的次数，以及去掉某信号后决策是否变化。这些是机制诊断，真正的收益归因仍需要在固定条件下比较调度策略。

冷启动是一个特殊情形。空历史下没有可比对象，输入新颖性可以按协议设为最大值，也可以设为缺失并另走启动规则；二者都不是数据自然告诉我们的答案。E11 v2 选择首题新颖性为一，所以事件臂可能在没有任何变化证据时就请求构建。这一请求来自初始化约定，应在记录里可辨认，不能称为成功发现了分布漂移。

\subsubsection{把报警、诊断与干预分开}

一个失败比例跳变可以要求检查，却不一定要求立即发布新记忆。跳变可能来自任务难度、服务暂态、验收器升级、工具故障，也可能来自已有记忆污染。对这些原因采取相同的“再写一条经验”，不但未必有效，还可能把暂态故障固化成长期规则。

一种工程设计是将响应分级：第一层记录异常并补齐观察；第二层构建候选但暂不发布；第三层在适用条件与回归检查通过后发布。另一个选择是先回退到已知版本，保留候选供离线分析。本书不声称这些分级一定提升质量，因为它们也增加等待和检查成本；它们的价值是让动作与证据强度相称，并使失败位置可以复查。

为减少在阈值附近反复触发，还可以研究冷却间隔、开启与关闭阈值分离、或要求持续越界等策略。这些都改变了可响应时间和遗漏风险，不能作为“纯实现优化”悄悄加入实验。尤其在短序列里，一个冷却期可能吞掉所有后续机会。应先把调度写成状态机，再列出边界输入逐项检查，最后决定是否值得付费评价。

\section{Gate 自己也是一次决策：阈值两侧的代价}
\label{sec:ch13-gate}

\begin{sloppypar}
触发 Gate 常被画成一条旁路：不值得写就绕过去，省下的就是写入那一次调用。
两篇一手工作把这条旁路的账报了出来，结论比“绕过去就省了”复杂：
\textbf{Gate 本身要算一次信号、要定一个阈值，还要承担漏触发}——
下面分别核对其token、准入延迟与召回等指标；这些不是两篇论文都完整报告的同一份生命周期账。
\end{sloppypar}

\begin{table}[htbp]
\centering
\caption{同步进化与 Gate 进化在同一批带噪对话上的对照（D-MEM Table 1，$\rho=0.75$，GPT-4o-mini）}
\label{tab:ch13-dmem}
\begin{small}
\begin{tabular}{>{\raggedright\arraybackslash}p{3.4cm}ccccc}
\toprule
做法 & 整体 F1 & 单跳 & 多跳 & 开放域 & 总 token \\
\midrule
A-MEM（同步演化） & \cited{0.336} & \cited{0.208} & \cited{0.365} & \cited{0.111} & \cited{1,648K} \\
RPE Gate（跳过 / 缓存 / 深演化） & \cited{0.369} & \cited{0.246} & \cited{0.412} & \cited{0.085} & \cited{319K}（\cited{$-80\%$}） \\
\bottomrule
\end{tabular}
\end{small}
\end{table}

\begin{sloppypar}
第一行是 D-MEM 论文采用的同步 A-MEM 比较条件。表中的 $\rho=0.75$ 保留原文设置标记：其 §2.2 先将新增噪声数与原始对话长度相比，紧接着又称最终序列中\cited{75\%}为噪声，两处的分母不一致\cite{dmem-2603}。若按前一种定义，最终噪声占比应为 $0.75/(1+0.75)=\calcnum{3/7}$，约\calcnum{42.86\%}。因此本表转述的是作者标记的配置，不能据此认定实际序列的噪声占比；复现时须另核对构造数据。
第二行采用分层路由：评判器将\textbf{惊讶程度}与\textbf{长期效用}组合成原文所称的 Agentic RPE；该指标不是通常的时序差分预测误差。
较低分数对应跳过或构建，较高分数触发深演化；冷启动另有强制构建规则，见第\ref{ch06}章。
结果不是一边倒：整体、单跳、多跳、对抗四项都涨了，
\textbf{但开放域那一项从 \cited{0.111} 掉到 \cited{0.085}}
。这组总体与分项变化说明存在取舍；它比较的是整个方法，不能把改善全部归因于触发时间。
总量上，token 从 \cited{1,648K} 降到 \cited{319K}，
本书按原表相除得到同步方案约为最终方案的\calcnum{5.17}倍，按差额除以原值计算下降约\calcnum{80.64\%}。倍率与降幅的分母不同。
公开归档的两份结果对记忆摄入费用的覆盖并不相同。同步 A-MEM 记录了摄入 token，最终 D-MEM 的各轮摄入计量均为零。两份结果的元数据均记录启用缓存。所指提交的代码在缓存命中时直接返回旧响应，跳过原始生成的 token 计量\footnote{公开结果取自作者仓库\url{https://github.com/london-and-tequila/dmem}的提交\texttt{68410d05}，计量路径按结果元数据所指的提交\texttt{e6f4043d}核对。}。因此，表中降幅适用于这两份运行记录。比较从空缓存完整建库的费用，还需要匹配初始化与缓存条件，并交代复用响应的原始生成费用。
\end{sloppypar}

\begin{sloppypar}
D-MEM 的 Table 3 比较了四个逐步修改的版本，不能当作只改变一个阈值的四档扫描。v1 使用线性 RPE，跳过率为\cited{9.6\%}，总量为\cited{489K} token；相对于前表同步方案，仍已减少约\calcnum{70.33\%}，不能称为“几乎没省”。v2 加强生命周期分类与乘性门，跳过率为\cited{68.5\%}，总体 F1 为\cited{0.291}。v3 放宽效用保护、取消实体提取开销；v4 在此基础上增加 BM25 混合检索和影子缓冲。v3 与 v4 的跳过率同为\cited{41.1\%}，token 却分别为\cited{531K}与\cited{319K}，比例约\calcnum{1.66}。这说明相同跳过率不能代表相同流程，也不能把二者差额全部算到触发规则上。原文称检索增强为“零成本”，其具体范围是没有额外大模型 token；本地索引、检索和存储仍需资源\cite{dmem-2603}。
\end{sloppypar}

\begin{sloppypar}
把准入写成显式判决的那一侧给出了阈值的形状。
一项把记忆准入拆成五个可解释因子（未来效用、事实置信度、语义新颖性、时效、内容类型先验）
并学出权重的工作报告（LoCoMo记忆准入测试候选 $N=\cited{225}$）：
F1 \cited{0.583} 对最强的对照 \cited{0.541}，相对 \cited{+7.8\%}，
延迟 \cited{2,644} 毫秒对 \cited{3,831} 毫秒\cite{amac-2603}。
这里的F1衡量记忆准入标签，延迟按候选计算；它们不是问答F1与整任务时延。更有用的是它的两张副表：
去掉\textbf{内容类型先验}这一项，$\Delta\mathrm{F1}$ 是 \cited{$-0.107$}
（\cited{0.583} 掉到 \cited{0.476}，正好退回到等权基线），
分别移除其余四个因子时，F1 各下降\cited{0.013}至\cited{0.028}，这些是单项消融，不能相加成为四项的联合贡献；
阈值敏感度表使用权重优化中的交叉验证折，给出 $\theta^{*}=\cited{0.55}$，不是另一次独立测试集搜索；
$\theta\le\cited{0.45}$ 时召回接近 \cited{1.000} 而精确率掉到 \cited{0.36--0.37}，
$\theta\ge\cited{0.65}$ 时漏掉 \cited{40\%} 以上的相关记忆，
而 \cited{0.50--0.60} 之间是一段平坦区。
\end{sloppypar}

\begin{pitfallbox}
\textbf{反例：把跳过率当成收益。}
跳过 \cited{68.5\%} 的输入听起来是省了三分之二，
原论文最终版与v2的整体F1分别为\cited{0.369}与\cited{0.291}，但二者还改变了其他组件。原作者把部分退化联系到基础上下文被过度过滤；这不是只改变跳过率所得到的独立因果效应。
\textbf{反例：把漏触发读成稳定。}
阈值定高之后，系统整段时间什么都不写；
从外面看，这与“系统已经很稳”完全同形。
\end{pitfallbox}

\begin{principlebox}
Gate可能节省生成、演化或后续检索费用，也可能因遗漏有用信息增加失败。节省发生在哪一步，应按实际路由和费用记录判断。
所以报告一个触发 Gate 时，要同时给出三件东西：
\textbf{（一）}Gate 前后的\textbf{成本}（token 与延迟各报一项，不能只报跳过率）；
\textbf{（二）}被跳过材料的覆盖、后续质量变化与不确定性；
\textbf{（三）}阈值选择过程与敏感性——如果不存在平坦区，也应如实报告，不能把平坦区视为方法成立的必要条件。
\end{principlebox}

\section{把一次更新拆成可检查的状态转换}
\label{sec:ch13-state-machine}

讨论“何时更新”时，最容易忽略的是更新并非一个不可分的按钮。任务完成、结果到达、信号计算、构建许可、候选产生、验证结束与正式发布，是不同的事件。把它们压成一个时间戳，就无法解释“已经触发却没有改变后续任务”到底发生在哪一层。本节给出本书的工程分解，用于理解 E11 v2 的记录方式，不将它当作所有系统唯一可行的架构。

\subsection{请求、尝试与发布是三个分母}

设第 $t$ 道任务完成后，调度器提出请求 $g_t\in\{0,1\}$。资源控制器可能因配额用尽而拒绝，所以实际构建尝试 $a_t$ 不必等于 $g_t$。构建返回后，候选又可能不满足准入要求，于是发布状态 $p_t$ 还可能为零。在这个简单设计里，三者满足
\[
  p_t\le a_t\le g_t.
\]
这只是事件之间的包含关系，不表示越靠左越接近“学会了”。发布只说明状态发生改变，还需要后续检验才能谈收益。若协议允许不经新构建直接恢复旧版本，则恢复应另设事件类型，不能硬塞进这条包含关系。

报告触发率时，分子必须说清。请求次数除以可观察任务数衡量信号有多频繁地要求动作；尝试次数除以请求次数表示请求转为实际构建的比例；发布次数除以尝试次数衡量候选通过了多少。末题抑制、配额拒绝与发送前的资源停止应分别记因，这个比例不能单独归因于预算。把三者都叫“更新成功率”，会把预算紧张解释为候选质量差，也会把格式准入解释为能力改善。分母为零时应报告不适用，不能为了图表整齐填成零成功率。

还要说明“可观察任务数”是否包括末题。E11 v2 仍计算末题信号，但明确不为这条已结束序列构建新经验。若把末题计入调度覆盖率，须保留“没有未来任务”的状态；若统计可行动机会，则可以排除末题，但分母需要同步缩小。两种口径都可描述同一过程，混用却会人为制造触发率差异。

以本章第~\ref{sec:ch13-trace-example}节的八题教学序列为例，周期臂在第\calcnum{2、4、6、8}题后提出请求，只有前三次获准构建。计入末题时，请求率为\calcnum{4/8}，获准数除以请求数为\calcnum{3/4}；若只统计可行动请求，后一个比例为\calcnum{3/3}。本例没有请求因构建配额被拒，差额来自末题规则。因此这些比例必须同时附上观察窗口和拒绝原因。控制器给出的构建许可也不证明调用已经发起，更不能代替发布数。

这套拆分也能解释预算应何时扣除。候选生成可能已经消耗模型资源，即使返回空文本或在发布前失败，也不能把尝试配额退回，再反复生成直到成功。否则上限约束的是成功次数，而失败成本可以无限增长。相反，若请求在任何调用之前就因资源不足被拒绝，则请求计数增加、调用计数不增加，两本记录不矛盾。E11 v2 在准许构建时占用尝试配额，在发送前再检查共同调用上限，并把任何进一步停止记录为不完整序列。

\subsection{有了候选，不代表持久状态已经改变}

候选可以先写入暂存区，经过格式、事实或回归检查后再发布。暂存区与正式记忆必须有明确边界：后续任务只能检索已发布版本，否则“尚未准入”的内容已经通过上下文影响行为，准入制度只剩表面形式。若研究问题就是让模型读取未验证候选，则应把它定义为另一种方法，不能在文中说只有验证过的经验进入任务。

E11 v2 的流程实现采用很窄的准入：单段、长度合规、非空且非完全重复。它没有独立判断经验是否正确，所以日志写“发布”，不写“验证为真”。构建输入包括来源任务描述、按初始文件和成功写入记录重建的可写源码，以及验收结果。通过只能说明源码满足该题验收；失败只说明存在不满足项。模型写出的原因解释，无论语气多确定，都不能由这个布尔判定直接推出。

例如某题修复后通过，候选说“所有循环上界都应加一”。这句话可以很短、没有重复、格式完全合法，也可能在下一题造成错误。来源通过与经验泛化正确，是两个不同层次的判断。因此 E11 首先考察这种固定更新器在不同调度下的表现；即使某种时机占优，结论仍依赖这个更新器。换成经过更严格事实核验的更新器，构建费与发布率都会改变，不能直接沿用原排序。

候选状态还应保留前后版本。只记“写了一条经验”，无法判断是新增、覆盖还是重复；只保存最终记忆，会抹掉中间任务实际看到的内容。本书在每题开始时记录记忆状态，在发布记录中保存新旧摘要值，并保留原始文本。摘要值有助于检查一致性，原始文本用于理解语义，两者不能互相替代。摘要值用于核对所编码的版本，适用条件仍须另查。

\subsection{更新发生在结果之后，收益发生在未来}

设任务 $t$ 使用状态 $M_t$，得到结果 $Y_t$，随后构建候选并发布为 $M_{t+1}$。这一顺序意味着当前更新不可能解释当前任务的成功：$Y_t$ 已经在发布之前形成。若图里把发布时刻与任务成绩放在同一列，必须注明作用方向。把当前题成功率回填给新经验，会让任何从成功题产生的经验都显得有效，这是时间标记造成的伪收益。

同样，一条经验是否来得及发挥作用，取决于发布延迟与后续机会。假设任务连续到达，而一次验证跨越了后面数题，那么这些题仍在使用旧状态；不能只按“构建请求已经发生”把它们划为更新后。最直接的记录是每个任务实际读取的版本号，而不是根据墙钟大致推断。版本号也有助于排查缓存：正式状态已更新但执行节点仍拿旧缓存时，时间戳显示更新发生了，输入记录却证明它没有进入决策。

对于有限实验，末题之后构建的经验只有成本，没有当前序列内的潜在收益。对于长期服务，末题只不过是人为截断，下一批任务仍可能使用它。因此成本截点不是技术细枝末节，而是估计对象的一部分。E11 小试选择只评价当前固定序列，所以关闭末尾构建；这不意味着部署系统永远应该这样做。若改成连续服务评估，应增加观察窗口和存量资产价值的范围说明，不能把窗口外收益凭空补到结果里。

\section{四笔账：更新、验证、存储与推理}
\label{sec:ch13-costs}

更新时机的选择之所以必须实测，是因为一次更新的代价分散在四本账上，
只算部分费用可能改变比较结论，是否反转仍需核算。

\begin{itemize}
  \item \textbf{更新账}：生成候选、改写关联条目与维护索引的实际工作。它可能只涉及一条经验，也可能扫描多条旧记忆；不能预设它一定最小。
  \item \textbf{验证账}：检查候选格式、事实、兼容性与回归表现。验证对象与测试范围决定代价，不能把论文讨论的整个记忆演化复杂度统称为验证复杂度。
  \item \textbf{存储与检索账}：持久保存、索引更新、查询与读取。库变大不意味着每次查询必定线性变贵；索引结构、缓存和检索上限都影响实际增长。
  \item \textbf{推理账}：取回材料注入上下文后带来的输入、输出与后续调用。相同检索条数可以有不同总长度，也可能改变解决任务需要的轮数。
\end{itemize}

这四笔账是过程归属，不是天然互斥的仪表。一次模型调用若同时判断候选并生成改写，可以在功能说明中涉及两项，但合计只能记一次。模块内的耗时也可能包含子模块耗时，不能把包含关系直接相加。对性能瓶颈的诊断可以采用多种视角，对总费用的核算却需要一份唯一的原始调用账。

JitRL 提供了另一个成本范围的例子。其 Table 9 报告 JitRL 约\cited{290}美元、WebRL 约\cited{9,900}美元；§5.7 与附录 R 说明，后者按训练硬件费用估算，前者主要是 API 使用\cite{jitrl-2601}。这支持特定设置下成本结构不同，不能把差额直接解释为“记忆不演化”这一个因素的因果效果。JitRL 仍通过累积经验影响后续策略，只是没有梯度更新。其附录 J.1、Table 13 的同基座、分离训练测试设置另报 JitRL 成本\cited{200}美元、平均成功率\cited{40.88}，WebRL 为\cited{9,900}美元、\cited{46.06}；这组数字也不能与 Table 9 的\cited{290}美元拼成同一条件。比较省钱时，必须让任务范围、基座、适应阶段和质量约束一同出现。

\subsection{预算分配与停止条件}
\label{sec:ch13-budget}

触发之后还有两个决策：这一次给多少预算，以及什么时候停。

给多少预算的问题在第 3 章已经用 episode 配额的形式解决过一次，这里的形式不同但结构相同：
\textbf{触发决定了“要不要花”，预算决定了“最多花多少”，停止条件决定了“什么时候不再花”}。
三者缺一，账就不闭合——只有触发没有停止条件的系统，在信号持续为真时会一直写下去。

资源停止条件应包含\textbf{可立即核对的账目上限}；质量下降、验证失败与安全条件也可以独立要求停止。它们不必互相替代。
质量的提升往往要等很久才看得出来（延迟反馈，见第\ref{sec:ch13-feedback-clock}节），
而“最近 $k$ 次更新里，有几次已经发布了可供下一题读取的新版本”可在发布时核对；下一题是否实际读取，仍需任务输入记录。
后者不直接等于收益，但它是一个诚实的中间量：
如果连续多次更新都没有改变下一次的起点，那么继续更新的理由需要从别处来。

\subsection{停止搜索与交付版本要分开决定}
\label{sec:ch13-stop-output}

一次候选是否准入、是否继续提出新候选、停止后交付哪个已保留版本，是三个不同决定。候选可能带来很小的验证增益，足以通过当前准入规则，却未必值得再付一轮搜索费用；拒绝当前候选，也不表示后续候选都没有价值。Yin等将停止演化与停止后的输出选择分别建模：前者累积候选增益不足的证据，后者估计这种变化从何时开始，再选择变化前的保留版本\cite{enough-enough-2610}。这使“花到哪一轮才决定停”与“最终交付哪一轮工件”不必相同。

这项方法还说明为什么被拒候选也值得保存。在原文的贪心准入设定中，当前保留版本只有在候选验证均分更高时才改变。如果只画保留版本的得分曲线，连续拒绝会形成平台，却看不见候选逐题改善、退化或持平的结构。作者使用同一验证题上候选与当前版本的配对结果积累停止证据，无论候选是否获准，都保留这些差异\cite{enough-enough-2610}。因此，“最近没有发布”可以是流程诊断，不能独自证明搜索机会已经耗尽。重复查看同一验证集也没有增加独立任务覆盖，最终质量仍需用途清楚的评价范围。

原文的“值得继续”由最低期望验证增益定义，其统计结论依赖相应的条件均值与独立性等假设；它没有将美元、等待时间和任务收益共同求解为最优调度。本书据此采用的是决策拆分：保留逐题候选对照，另设停止规则与输出选择规则，并保存可能被返回的版本。若还需要额外评价，费用仍进入总账；硬资源上限也始终有效。停止时返回较早版本，不会撤销之后已经发生的调用、写入或外部副作用。

\subsection{共同上限、实际费用与单位资源收益}
\label{sec:ch13-resource-estimand}

E11 v2 给每个臂的整段序列相同的调用上限，同时保留每题的调用上限。这样能防止某臂无约束地多做工作，但不等于三臂实际消耗相同。一个策略提前结束任务，另一个策略多读几轮错误输出，输入长度和调用数都会变化。让所有臂在最后补空调用以“花满预算”，只会破坏资源效率问题的含义。

\subsubsection{先定义要比较哪一种效率}

至少有三种不同问题。第一是在共同资源上限下，哪个策略达到更高质量；第二是在达到预定质量要求时，哪个策略实际资源更少；第三是质量变化与资源变化的某个标量折衷。第一种可以报告受限成绩和实际账，第二种需要明确质量门槛及其不确定性，第三种必须在看结果前说明折衷系数。只写“单位资源收益更高”，不能替代这些定义。

一个教学反例足以说明比值的局限：方案甲答对\calcnum{8}题、消耗\calcnum{10}个资源单位，方案乙答对\calcnum{6}题、消耗\calcnum{5}个单位。乙的成功数除以费用更高，但若应用要求至少答对七题，乙仍不合格。若资源单位是调用次数，相同次数也未必同费用；若单位是美元，还要给出对应时间的价格和输入输出分项。比值可以描述结果，却不能自行决定接受条件。

为了减少这类歧义，可以先展示质量与资源的二维位置。一个策略在质量不差且资源更少时，至少在这两个观测维度上更有吸引力；一个策略质量更高、资源也更高时，需要读者的约束来判断。带噪估计还不能只比较点的位置，尤其不能把“不显著变差”当作满足质量不劣要求。第\ref{ch03}章的非劣边界与缺失结果范围在这里同样适用。

E11 v2 的三策略真实模型比较尚未运行。当前分析器按独立序列单元计算三对质量差与前后优势变化，并保留资源账；只有离线构造数据时，这些输出用于检查计算流程，不能估计真实策略的效率排序。

\subsubsection{一条显式的节约条件}

在一个简化教学模型里，设有 $N>0$ 次可更新机会，每次深更新成本固定为 $c_u>0$，事件信号每次成本为 $c_s\ge0$，事件策略还承担一次固定设置费 $F\ge0$，实际触发比例为 $r\in[0,1]$。忽略策略对后续推理的影响时，两种成本为
\[
 C_{\mathrm{each}}=Nc_u,\qquad
 C_{\mathrm{event}}=Nc_s+rNc_u+F.
\]
事件策略严格省资源需要
\[
 r<1-\frac{c_s}{c_u}-\frac{F}{Nc_u}.
\]
右边可能为负，此时在这些假设下，即使一次深更新也不做，监测与设置费仍已超过每任务更新。右边也会随运行长度变化，所以短期试跑与长期部署未必出现同一个回本位置。这是代数条件，不是已经测出的阈值。

例如教学设定 $N=\calcnum{100}$、$c_u=\calcnum{10}$、$c_s=\calcnum{1}$、$F=\calcnum{50}$ 时，需有 $r<\calcnum{0.85}$ 才严格省资源。这里没有货币单位，也没有声称真实模型符合固定成本假设。若经验令后续每题多读一段长文本，就应额外加入推理增量；若减少了反复排错调用，该增量也可能为负。这个简式的用途，是帮助识别还缺哪几笔账，而不是替代实测。

当一次触发有不同深度时，单个 $r$ 更不够。跳过、只构建、全演化可以分别有不同概率和条件成本；总体费用应按各路径实际次数与实际消耗相加。D-MEM 的分层路由正提示了这个区别\cite{dmem-2603}。不能把所有非跳过事件都按同一个“写一次”收费，也不能因为最便宜路径频繁出现就认定最贵路径不会主导总账。

\subsubsection{调用、请求尝试与费用为什么不能互换}

一次逻辑上的模型调用可能经历多次网络尝试。超时并不总能说明服务端没有执行；重试成功后返回的 usage 通常只说明这次可见响应的使用量。把最终成功响应的费用当作全部费用，会漏掉未知部分；把每次超时都按一次完整成功收费，又会加入未经证实的数字。最诚实的账是已知用量与未知尝试同时保留。

新版 E11 关闭后端内部隐式重试，由外层记录每一次传输调用。发送前写入待完成条目，结束后填入响应或错误类别。若进程在等待时退出，磁盘上仍有待完成项；它不是免费调用，也不是可确认为一次服务端执行的证据。通过这种状态记录，后续可以与提供方账单核对，而不必从完整任务数反推所有费用。

调用上限可以在发送前严格检查，因为下一次发送会不会增加一，是执行器能够控制的。输出 token 上限若只是收到回复后累计，就只能阻止下一次调用，不能追回已经生成的内容；墙钟同理，正在等待的请求或正在执行的本地验证可能跨过阈值。协议应写明这一差别。E11 小试给出的请求次数上限不构成金额硬封顶，也没有把观察到的已知 token 合计称为完整账单。

完整原始记录还防止另一种错误：把 episode 的汇总费用、调用层的费用和序列的费用各加一次。它们很可能是同一笔消耗的三个视图。新版以调用层为唯一合计来源，episode 汇总用于诊断和核对，不再累加到总量。本地监测耗时单列，整段墙钟已经包含其中一部分；描述分项时必须指出包含关系，不能把嵌套时间相加成一个更大的“真实总耗时”。

\section{同一条判定序列上的三种调度}
\label{sec:ch13-trace-example}

用一个有限的教学序列，可以看清时机与配额的相互作用。设\calcnum{8}题的判定依次为“通过、通过、通过、通过、失败、失败、失败、通过”，任务描述完全相同，三臂都允许最多\calcnum{4}次经验构建。这里把判定作为外部输入固定下来，只演示调度器；真实运行中，记忆会影响判定，三臂未必得到同一序列。

每任务臂在第\calcnum{1、2、3、4}题后获准构建并用完配额；此后仍逐题提出请求，但后续非末题请求被配额拒绝。周期臂在第\calcnum{2、4、6}题后获准构建。事件臂获准构建的位置为第\calcnum{1、5、6、7}题后：第一题是冷启动新颖性，第五题的短窗失败率差恰好达到阈值，第六、七题又满足连续失败或窗口变化。第八题虽同时满足周期与事件臂的请求条件，两臂都由末题规则阻止构建；请求仍保留在日志中。这些位置描述的是构建许可，不是全部请求或成功发布。决策由离线控制器逐项算出，未调用模型。

这个例子揭示了两种容易误读的现象。第一，事件臂并没有节省构建次数，它与每任务臂都用了四次，只是把更多机会留到后半段。第二，第五题只有一次失败，事件臂已经请求更新，所以不能把它的这次动作解释为“等到连续两次失败才出手”。逻辑或中的窗口信号先越界了。若只保存最终触发位而不保存每路信号，事后的文字解释很容易把设计说错。

配额也改变了“每任务”的实际含义。在无配额的教科书描述中，它每题都更新；在共同配额下，它更准确地说是“每题都请求，按先到先用消费”。这种先用完策略可能吃亏，但不能在看到结果后偷偷给它重置配额来修正。若研究者希望比较的是均匀分配的每任务策略，就需要为每次构建分配更小预算、改变准入或调整实验问题，形成另一份协议。

同理，周期的相位也重要。每两题之后更新与从第一题开始每隔两题更新，并非同一个策略；漂移位置若恰好落在某一相位附近，短序列结果会受其影响。确认性设计可以预先安排多种切换位置与相位，或者把它们作为设计因素；单次工程试跑只能把采用的相位准确报出，不能泛化为所有周期策略。

即便事件臂在后半段分数更高，也仍要检查更新内容与额外调用。它可能从失败中得到更有诊断力的材料，也可能只是多花了做题轮数；它在第一题构建的经验还可能一路影响后面的触发输入。时机、内容与状态是闭环关系，不能把一张触发位置图直接翻译成“发现漂移所以提高成功率”的因果链。

\section{条件变化：换任务族、换模型、换反馈器}
\label{sec:ch13-conditions}

触发阈值是在具体观察范围内选择的。换任务后需要重新检验校准与资源表现，但不能先在新的测试结果上任意调好，再把它叫作原策略的迁移。

换任务族时应重新检查\textbf{失败率基准}。
分布漂移这一路信号比的是“最近窗口的失败率与此前的差”，
它的数值分布依赖该族的基准失败率和错误相关性，虽然量纲始终是比例差：
在一个基线失败率为 \measured{0.25} 的训练过程里（E01 格 4，\measured{8} 个 episode），
旧脚本将七个暖机位置填零后得到均值\measured{0.0625}，真正形成两窗口的比较仅有一次；
换到一个基线失败率接近 0.5 或接近 0 的族，同一个阈值的触发概率和决策代价可能不同。
因此应报告任务族组成并设置稳定分布对照，或在可比族内检查变化。按族归一化需要足够样本与预定方法，并非自动消除混淆。

换模型时还应检查\textbf{错误集中}这一路。
连续失败不要求每次错误原因相同；它只记录一串失败。换模型会改变边际失败率、错误相关性及反馈可见性，这几种变化都可能影响触发频率。一个总体更强的模型仍可能在某个新子域集中失败；“更不稳定”也不必然意味着连续失败更多，须检查错误出现的时间结构。

换反馈器时，还要检查反馈何时到达。旧的真实结果仍可用于延迟观察，只是其时间含义改变；不必把所有信号都改成代理量。若要使用置信度、评委评分等提前触发，应与尚未到达的真实验收分别记录。

\begin{sloppypar}
还有一维是\emph{监测本身的费用}，它最常被漏记：
三个信号都要在每次 episode 之后算一遍，虽然本书的三个信号都不调用模型
（\texttt{error\_run} 与 \texttt{drift} 只用已发生的验收结果，\texttt{novelty} 只做字符二元组重叠），
但换成需要一次模型调用的新颖性度量之后，
“每次都监测”就从本地计算扩展到模型调用，其量级是否接近每任务更新取决于提示、输出和执行频率。
\textbf{监测费上升会压缩事件触发的净收益，但是否逆转比较仍取决于节省量与任务表现}——
这也正是 E11 把监测开销写进账本的原因（§\ref{sec:ch13-e11}）。
\end{sloppypar}

\begin{sloppypar}
最后一维是\emph{误触发与漏触发的代价不对称}。漏触发可能让错误状态持续，也可能只错过一条价值很小的经验；误触发除了多花一次构建费，还可能把错误内容发布到后续任务，产生持续污染。因此不能把前者统一当作随时间增长的损失、后者统一当作固定常数。阈值取向需要结合更新可逆性、验证能力、未来复用次数与两类错误的估计损失。若这些量尚未测清，保留候选而推迟发布，往往比一律放宽阈值更便于诊断。
\end{sloppypar}

\section{迟到的反馈如何进入更新循环}
\label{sec:ch13-feedback-clock}
\label{sec:ch13-delayed}

快速在线反馈只是本章小型编程实验的条件，部署系统可能几小时后才得到验收，也可能只收到部分用户的反馈。此时可以使用已经到达的真实结果，也可以提前使用代理信号，两者回答的问题不同。已经到达的结果描述过去执行状态；代理信号描述当前可见线索，都不能自动还原尚未观察到的实际收益。

E02 已完成有限任务上的反馈策略比较，见第\ref{ch04}章；它不是这里所有代理信号的校准实验，也没有证明部署延迟可以忽略。需要为当前信号单列校准记录：当时可见的输入、使用的代理分数、最终到达的真实判定、等待时长与是否触发。不能用后来才知道的结果回写当时的输入，再声称触发器能提前发现问题。

在线服务至少有两只时钟：任务发生的时钟与结果到达的时钟。一个昨天执行的任务可能今天才被判失败；今天刚结束的任务可能尚无结果。如果直接按收到反馈的顺序计算“连续失败”，得到的是反馈到达流上的连续性，不一定是执行任务流上的连续性。两者都可以定义，但不能在解释时互换。

一种做法是按执行序号保存待定位置，等相应结果到达后更新已有记录。这可以保留原任务顺序，却会让最近窗口包含缺失项，必须预先规定是否等待完整窗口、使用已到达子集，或只运行不依赖真实判定的信号。另一种做法是按反馈到达顺序运行检测器，但把它明确称为到达流检测。不能把未到达的结果填成成功，也不能填成失败；前者压低报警，后者让延迟本身制造报警。

反馈缺失还可能是选择性的。用户更愿意投诉严重错误，而普通成功常常没有明确确认；某类复杂任务的验收又可能更慢。只在有反馈的样本上计算成功率，会混合任务质量与反馈机制。记录反馈覆盖率、等待分布与任务类型，至少能让这种混合可见。用代理分数补齐时，应保留哪些是真实判定、哪些是预测，不能合成一列无标识的“结果”。

对于触发决策，应保留决策时刻的可见信息快照。真实判定后来到达，可以用于分析该触发是否合理，却不能反过来改写当时的信号。离线回放若直接读取一份已经补齐的完整表，也可能无意使用未来信息。最简单的检查是为每个反馈附到达时间，回放到某时刻时只释放已经到达的记录；否则所谓提前检测，只是离线观察者知道了在线系统尚不知道的答案。

延迟还会影响在途更新。假设同一个旧版本产生的多条失败反馈陆续到达，第一条已经触发候选构建，后几条是否再次触发？若重复构建同一修复，资源会浪费；若一律忽略后续反馈，又可能错过不同原因。可以把触发事件关联到来源版本与故障假设，合并同一问题的证据，同时保留新增原因。合并规则是待验证的工程策略，必须记录被合并的请求数量，而不是让日志看起来这些请求从未出现。

在反馈较慢时，暂存与发布的分离更有用。系统可以提前构建候选，等待回归结果后再发布；也可以先在隔离任务上评估，避免把未完成验证的内容暴露给主要任务流。这会增加等待和评估费用，适不适合取决于错误传播代价。对于可以快速撤销、影响很小的改动，较短等待可能合理；对于跨任务传播广的资产，应把恢复范围和回滚成本纳入时机选择。第\ref{ch14}章将继续讨论这种持久改变的负面影响。

\section{从工程试跑到可以裁判的时机实验}
\label{sec:ch13-evidence-ladder}

E11 v2以共享调用账约束做题与经验构建，序列目录拒绝覆盖。装置检查覆盖时序、路径、发布与费用累计，详见附录~\ref{app:lab}；三策略的能力与资源比较尚未运行。

后续流程小试可先固定有限任务流、三种时机和构建机会，检查一次完整序列；相应草案与资源上限保存在\texttt{experiments/protocols/}。通过装置预检后，仍要冻结代码、任务、信息可见时间与实际后端身份，再取得单独的真实运行记录。

工程小试首先回答流程能否完整结束、预算是否如定义生效、候选是否真正进入下一题、失败和中断是否留下可复核记录。它还可暴露响应格式、服务稳定性或单题预算不足等现实问题。如果出现这些问题，应保留本次结果与费用，修订为新版本后再运行。把修前部分与修后部分拼成一次实验，会让三臂不再共享同一实现。

工程检查通过之后，才适合安排更有判决力的设计。P5 包含至少两个不同主张：某种时机在指定质量与资源条件下优于另一种，以及优势在漂移条件下扩大。前者需要足够的独立配对任务流与资源定义；后者还需要把漂移影响与任务难度、相位及切换方向区分。只有一个族切换、一个固定服务时段和一次重复，即使点差很大，也无法自动提供这些证据。

确认性设计可以预先安排稳定序列与变化序列，让三种策略分别面对同一个序列单元；还应考虑切换方向和切换时刻。若两种方向来自同一组题目，它们也未必是完全独立单位，重采样结构要反映共享来源。增加同一条固定序列上的模型重复，可以观察服务与生成波动，却不会自动增加任务族的覆盖范围。每种重复到底代表什么，需在分析之前写明。

三策略排序还要求直接检查相邻对比。事件优于每任务并不推出事件优于周期，也不推出周期优于每任务。若先看结果再只留下最好看的两臂，就把一个三策略问题变成了选择后的二策略展示。新版分析器同时输出三个预定对比，并把单簇的区间留空；后续若进行多重主张的正式检验，应另行定义控制方法与足够样本，不能把探索性区间自动当作联合结论。

最后，成绩差值的前后变化也有解释范围。换族后事件臂相对优势扩大，可能与新族更适合生成经验有关，也可能与基线在该族更弱有关；这并不单凭一个差值之差就证明事件检测器识别了“漂移本身”。需要结合信号何时越界、什么候选被构建、哪些任务实际读取了它，以及稳定对照中的对应行为。统计摘要和过程证据相互补充，任何一方都不能替另一方完成论证。


\subsection{阈值怎样选择，才能保留测试的判别力}

一个调度器通常包含不止一个可调量：失败连续长度、窗口宽度、差值阈值、新颖性阈值、冷启动规则与写入配额。即使每项只有少量候选，组合起来也可能形成很多策略。如果在同一批最终任务上不断试，直到某个组合看起来优于周期策略，报告的差异就混入了选择过程。固定随机标签或给每个组合相同调用上限，都不会消除这种选择偏差。

一种可审查的流程是把阈值开发与最终比较分开。开发阶段使用明确的训练或选择序列，记录考察过的组合、选择指标与总费用；选定后冻结策略，再在未参与选择的任务流上比较。若序列来自同一批模板的轻微变体，拆分也不能只看文件名不同，应说明任务族和生成来源是否共享。冻结后遇到缺陷可以修，但修改将形成新的版本与新的评价范围，不能继续把已经看过的测试流称为首次检验。

选择指标本身也要匹配主张。若用成功率选阈值，最后主张的是单位资源收益，则开发过程可能偏向更多调用；若只用触发次数选阈值，又可能选出几乎不更新的策略。可以预先设置质量门槛并在合格候选里比较费用，也可以列出质量与资源的候选集合，再用明确的应用偏好选择。无论采用哪种形式，所有被试过而未被选中的候选仍属于开发账，不能只留下赢家那次运行费。

A-MAC 的阈值敏感度表提供了一个具体的阅读提醒：其最优阈值是在用于权重优化的交叉验证折上选择，表中的验证 F1 与独立测试结果并非同一列\cite{amac-2603}。这使我们可以看见选择范围，却不意味着把该阈值搬到另一种任务或另一种记忆构造器后仍最优。所谓平坦区也应绑定对应数据与指标；它不是一个对所有应用都有效的安全区间。

运行期自动调阈值则又多了一层状态。阈值本身开始随反馈变化，比较对象已经从固定触发器变成自适应调度策略。此时应记录每个时点使用的阈值、更新依据和可用预算，并说明调整是否读取了当前评价任务的反馈。把这种调整称为系统功能完全可以，但不能一面让它在线挑参数，一面把对照限制为固定参数，却把差异仅归于触发信号更好。

\subsection{回放检查能回答什么，不能回答什么}

把已保存的判定轨迹喂给几个调度器，是低成本的开发手段。它能够检查冷启动、窗口暖机、周期相位、配额用尽和末题处理，比较在同一观察序列上会提出多少请求。第\ref{sec:ch13-trace-example}节采用的正是这种方式。对于控制器实现，它非常有用；对于新策略的未来成功率，它却存在一个根本限制：日志中的结果来自原来的状态与动作，未必是替代调度策略会产生的结果。

例如原策略在第三题后写入了错误经验，后面出现连续失败。离线回放显示另一策略不会在第三题触发，但仍把原日志里的后续失败喂给它，就把本来可能不发生的错误当成了它的观察。反过来，原策略早期构建的有效经验可能让后续任务成功，替代策略若没有构建这条经验，也不应免费继承这些成功。闭环中动作会改变未来观察，静态回放无法自动给出反事实轨迹。

因此本书把回放结果命名为调度诊断，并明确固定了观察输入。要评价完整策略，需要让各臂各自执行后续任务，保留实际读到的记忆与新反馈，同时确保初始任务流可比较。更复杂的离线策略评价需要额外的行为记录与识别假设，不能从一列成功失败结果直接推出。没有这些条件时，诚实地说“能比较请求位置，不能估计替代策略的成绩”比画一条虚假的收益曲线更有用。

在部署前，还可以准备一组不涉及模型能力的反例输入：重复反馈、反馈乱序、长时间没有结果、窗口尚未形成、预算恰好耗尽、请求返回空内容、进程在发送后中断。每一种输入都要求一个可解释的状态，而不是笼统返回失败。新的 E11 流程检查覆盖其中一部分，但没有覆盖全部生产并发情形。实验台的单写入者和串行任务假设需要保留在说明中，不能把离线通过升级为分布式运行可靠性保证。

上线后的审计入口也应尽量直接。读者或维护者应能从某题成绩找到它实际使用的版本，从该版本找到来源候选和发布记录，再从候选找到原始请求与可见反馈。沿着这条链，才可能区分“选错时机”“学到错误内容”“正确内容没有被检索”和“验收器改变”几类原因。把所有问题都归咎于模型不稳定，会让下一轮调参缺乏方向；把每一次成功都归功于触发器，也同样超出了记录能够支持的范围。

% figures/ch13-trigger.tex —— 第 13 章核心图（全书引用入口）
% 用法：在 book/chapters/ch13.tex 中 % figures/ch13-trigger.tex —— 第 13 章核心图（全书引用入口）
% 用法：在 book/chapters/ch13.tex 中 % figures/ch13-trigger.tex —— 第 13 章核心图（全书引用入口）
% 用法：在 book/chapters/ch13.tex 中 \input{figures/ch13-trigger}
% 图主体在 ch13-trigger-body.tex，便于单独编译与复用。
\begin{figure}[htbp]
\centering
\input{figures/ch13-trigger-body}
\caption{固定判定输入下的调度教学例：配额为\calcnum{4}次、周期为\calcnum{2}题、末题不构建。事件臂在第\calcnum{1、5、6、7}题后获准构建；首题来自冷启动，第五题的短窗变化恰好达到阈值。末题虽有请求，仍不构建。构建标记表示获准尝试，不代表发布成功或能力收益。实际模型运行的三臂判定可能不同。}
\label{fig:ch13-trigger}
\end{figure}

% 图主体在 ch13-trigger-body.tex，便于单独编译与复用。
\begin{figure}[htbp]
\centering
% Teaching schedule computed in calculations/ch13_calculations.py; no model outcomes.
\begin{tikzpicture}[x=1cm,y=1cm,font=\scriptsize,
  cell/.style={draw,rounded corners=1pt,minimum width=0.96cm,minimum height=0.5cm,inner sep=1pt},
  build/.style={cell,fill=blue!12},
  stop/.style={cell,fill=gray!12},
  note/.style={align=left,text width=13cm,font=\footnotesize}]
  \foreach \i in {1,...,8} {
    \node at (2.1+1.35*\i,0) {题\i};
  }
  \node[anchor=east] at (2.8,-0.65) {固定判定};
  \foreach \i/\v in {1/通,2/通,3/通,4/通,5/败,6/败,7/败,8/通} {
    \node[cell] at (2.1+1.35*\i,-0.65) {\v};
  }
  \node[anchor=east] at (2.8,-1.45) {每任务};
  \node[anchor=east] at (2.8,-2.2) {周期};
  \node[anchor=east] at (2.8,-2.95) {事件};
  \foreach \i in {1,2,3,4} {\node[build] at (2.1+1.35*\i,-1.45) {构建};}
  \foreach \i in {5,6,7} {\node[stop] at (2.1+1.35*\i,-1.45) {额满};}
  \foreach \i in {2,4,6} {\node[build] at (2.1+1.35*\i,-2.2) {构建};}
  \foreach \i in {1,3,5,7} {\node[cell] at (2.1+1.35*\i,-2.2) {等待};}
  \foreach \i in {1,5,6,7} {\node[build] at (2.1+1.35*\i,-2.95) {构建};}
  \foreach \i in {2,3,4} {\node[cell] at (2.1+1.35*\i,-2.95) {等待};}
  \foreach \y in {-1.45,-2.2,-2.95} {\node[stop] at (12.9,\y) {末题};}
  \node[note,anchor=north west] at (0.5,-3.55) {三层事件分别记录：提出请求 $\rightarrow$ 配额准许构建 $\rightarrow$ 候选准入与发布。\\共同上限约束机会；实际调用、重试、已知用量与未知费用另行记账。};
\end{tikzpicture}

\caption{固定判定输入下的调度教学例：配额为\calcnum{4}次、周期为\calcnum{2}题、末题不构建。事件臂在第\calcnum{1、5、6、7}题后获准构建；首题来自冷启动，第五题的短窗变化恰好达到阈值。末题虽有请求，仍不构建。构建标记表示获准尝试，不代表发布成功或能力收益。实际模型运行的三臂判定可能不同。}
\label{fig:ch13-trigger}
\end{figure}

% 图主体在 ch13-trigger-body.tex，便于单独编译与复用。
\begin{figure}[htbp]
\centering
% Teaching schedule computed in calculations/ch13_calculations.py; no model outcomes.
\begin{tikzpicture}[x=1cm,y=1cm,font=\scriptsize,
  cell/.style={draw,rounded corners=1pt,minimum width=0.96cm,minimum height=0.5cm,inner sep=1pt},
  build/.style={cell,fill=blue!12},
  stop/.style={cell,fill=gray!12},
  note/.style={align=left,text width=13cm,font=\footnotesize}]
  \foreach \i in {1,...,8} {
    \node at (2.1+1.35*\i,0) {题\i};
  }
  \node[anchor=east] at (2.8,-0.65) {固定判定};
  \foreach \i/\v in {1/通,2/通,3/通,4/通,5/败,6/败,7/败,8/通} {
    \node[cell] at (2.1+1.35*\i,-0.65) {\v};
  }
  \node[anchor=east] at (2.8,-1.45) {每任务};
  \node[anchor=east] at (2.8,-2.2) {周期};
  \node[anchor=east] at (2.8,-2.95) {事件};
  \foreach \i in {1,2,3,4} {\node[build] at (2.1+1.35*\i,-1.45) {构建};}
  \foreach \i in {5,6,7} {\node[stop] at (2.1+1.35*\i,-1.45) {额满};}
  \foreach \i in {2,4,6} {\node[build] at (2.1+1.35*\i,-2.2) {构建};}
  \foreach \i in {1,3,5,7} {\node[cell] at (2.1+1.35*\i,-2.2) {等待};}
  \foreach \i in {1,5,6,7} {\node[build] at (2.1+1.35*\i,-2.95) {构建};}
  \foreach \i in {2,3,4} {\node[cell] at (2.1+1.35*\i,-2.95) {等待};}
  \foreach \y in {-1.45,-2.2,-2.95} {\node[stop] at (12.9,\y) {末题};}
  \node[note,anchor=north west] at (0.5,-3.55) {三层事件分别记录：提出请求 $\rightarrow$ 配额准许构建 $\rightarrow$ 候选准入与发布。\\共同上限约束机会；实际调用、重试、已知用量与未知费用另行记账。};
\end{tikzpicture}

\caption{固定判定输入下的调度教学例：配额为\calcnum{4}次、周期为\calcnum{2}题、末题不构建。事件臂在第\calcnum{1、5、6、7}题后获准构建；首题来自冷启动，第五题的短窗变化恰好达到阈值。末题虽有请求，仍不构建。构建标记表示获准尝试，不代表发布成功或能力收益。实际模型运行的三臂判定可能不同。}
\label{fig:ch13-trigger}
\end{figure}


% ---------- 6. 选择原则 ----------
\begin{principlebox}
\textbf{先报触发率，再谈时机好坏。}
周期调度、配额限制和末题不更新也会减少次数；
不报触发率，“事件触发更省”就只是个说法。

\textbf{信号要有分辨率，不只是要有定义。}
本书历史文本的新颖性均值为\measured{0.9546}，按\calcnum{0.5}阈值回算全部放行。这个观察说明该数据与阈值组合没有筛选作用，不能证明信号处处无用。应同时报告分布、阈值及是否实际部署。

\textbf{把四笔账分开记。}
更新、验证、存储、推理各自的量级不同，随规模增长的速度也不同；
只报总量会掩盖分项，只有分项又无法判断净成本；应保留不重复计数的总账与分项视图。

\textbf{资源上限和质量停止条件分别写明。}
发送前可以检查执行器能控制的调用次数上限；token与墙钟采用软停止阈值时，还需说明一次回复或执行跨过阈值的范围。未知费用和迟到的质量反馈仍需保留，并可能要求停止发布或回滚。

\textbf{把不请求、被拒绝、执行失败和发布失败分开。}
触发率为零不能证明稳定，触发率很高也不能证明有效更新。还要检查反馈覆盖率、候选状态和后续任务质量。
\end{principlebox}

% ---------- 7. 反例 ----------
\begin{pitfallbox}
\textbf{反例一：把离线回算写成已部署策略。}
旧运行的\measured{8}条经验按\calcnum{0.5}阈值回算全部放行；原运行没有执行这个新颖性门。应报告回算身份与固定输入，不声称实际节约或闭环收益。

\textbf{反例二：把“少更新”当成“会挑时机”。}
触发率下降可能来自阈值、输入分布、配额或检测能力变化，单个频率无法区分原因。
JitRL 的无梯度方案与训练方案费用结构不同\cite{jitrl-2601}；它依然积累经验，不能把其成本表解释成一个已经学会更新时机的消融。

\textbf{反例三：只记更新账。}
只记录经验生成会漏掉事实核验、关联改写、索引维护与后续推理。D-MEM 对特定演化路径的复杂度讨论\cite{dmem-2603}不能跨论文充当所有系统验证账的统一增长率；应在同一设计里逐项测量。
\end{pitfallbox}

\section*{习题}
\begin{enumerate}[label=13-\arabic*.,ref=13-\arabic*]
  \item\label{ex:ch13-1} 取你熟悉的一个系统，指出它属于表~\ref{tab:ch13-timing} 的哪一行；
        估算它的更新次数与任务数之比。请分别定义请求率、尝试率与发布率，说明哪些分母在该系统中相同。
  \item\label{ex:ch13-2} 本书\measured{8}条已存记录的新颖性都达到\calcnum{0.5}阈值，但分数并不相同。
        换一个相似度（例如先抽关键词再比），你会怎样区分“分数有变化”“能识别语义新信息”与“值得发布”？说明需要哪些标注、怎样选择阈值，以及如何保留独立评价。
  \item\label{ex:ch13-3} 四笔账各在什么实现条件下可能超线性增长？在你的系统里怎样测量，并避免把累计费用当单次费用？
  \item\label{ex:ch13-4} 一个事件触发系统近三天的触发率是 0，同期任务成功率也稳定。
        列出三种可能的解释，并说明你会加记哪个量来区分它们。
\end{enumerate}

\section*{常见误解}
\begin{itemize}
  \item \textbf{“事件触发一定比每任务更省”。}
        总费用还取决于监测、候选长度、失败重试和后续任务行为。即使更新次数相同，更新落点与内容也可能不同，不能直接称两者等价。
  \item \textbf{“新颖性低就说明这条经验没用”。}
        新颖性衡量的是与已有经验的重复程度，不是有用程度；
        一条高度重复的经验可能恰恰是最该被写进去的。
  \item \textbf{“延迟反馈下的代理信号可以先用了再说”。}
        代理信号与真实收益的相关性本身需要被测；
        未校准就长期使用，等于把触发权交给一个未经验证的量。
\end{itemize}

\section*{进一步阅读}
\begin{itemize}
  \item \cite{enough-enough-2610}：将何时停止搜索与停止后返回哪个保留版本分开，使用含被拒候选的逐题配对结果；其验证增益门槛不是费用最优解。
  \item \cite{dmem-2603}，讨论同步演化的规模成本与 RPE 分层路由。
        其带噪对话设置可为触发比较提供设计参照，实际噪声比例与构造流程须先核对。
  \item \cite{amac-2603}：把准入做成显式判决，
        Table 2中移除类型先验造成最大点值下降，支持该配置下的条件贡献，不能推广为普遍重要性排序。
  \item \cite{jitrl-2601}：无梯度经验积累与参数训练采用不同费用口径（\cited{290} 对 \cited{9,900} 美元），
        可用于辨析省下的资源来自哪些过程。
  \item \cite{agemem-2601}：阶段边界如何切分上下文与记忆，
        这一阶段安排与固定周期发布不是同一个机制。
\end{itemize}
